\documentclass[12pt]{article}
\usepackage{amsmath, amssymb, amsthm, geometry}
\geometry{a4paper, margin=2.5cm}

% XeLaTeX + русский язык
\usepackage{polyglossia}
\setdefaultlanguage{russian}
\setotherlanguage{english}
\usepackage{fontspec}
\setmainfont{Times New Roman}

\usepackage{enumitem}
\usepackage{xcolor}
\usepackage{tcolorbox}

\newtcolorbox{critbox}{colback=red!5, colframe=red!70!black, title=Критическая ошибка, fonttitle=\bfseries}
\newtcolorbox{modbox}{colback=orange!5, colframe=orange!70!black, title=Существенное замечание, fonttitle=\bfseries}
\newtcolorbox{minbox}{colback=blue!5, colframe=blue!60!black, title=Мелкое замечание / сомнение, fonttitle=\bfseries}

\newcommand{\R}{\mathbb{R}}
\DeclareMathOperator{\sign}{sign}

\title{Валидация доказательства\\«О невозможности сильного понижения размерности\\ в псевдоевклидовом аналоге леммы Джонсона--Линденштраусса»}
\author{Отчёт о проверке}
\date{\today}

\begin{document}
\maketitle

\section*{Краткое резюме}

Ядро работы — комбинаторно-геометрическая оценка (Теорема~1, раздел~3) — \textbf{корректно}: оценка числа реализуемых матриц через число знаковых условий полиномов (Милнор--Том / Бэзу--Поллак--Рой) проведена верно, и вывод о существовании нереализуемой матрицы при $p+q<Cn$ обоснован.

Однако \textbf{применение к лемме Джонсона--Линденштраусса (раздел~4) содержит критический дефект в самом определении $\varepsilon$-изометрии}: неравенство~(2) в псевдоевклидовом случае \emph{неудовлетворимо} для пар с отрицательным квадратом расстояния, а ключевой шаг «$f$ сохраняет ненулевые расстояния» обоснован только для положительных квадратов расстояний. Кроме того, есть переусиленное утверждение о «единой плохой матрице», неверная библиографическая ссылка и ряд концептуальных сомнений. Ниже — по пунктам.

\section{Критические ошибки}

\begin{critbox}
\textbf{1. Неравенство~(2) неудовлетворимо для отрицательных квадратов расстояний.}

Определение $\varepsilon$-изометрии дано как
\[
(1-\varepsilon)\|x-y\|^2 \le \|f(x)-f(y)\|^2 \le (1+\varepsilon)\|x-y\|^2. \tag{2}
\]
В псевдоевклидовом пространстве $\|x-y\|^2$ \emph{может быть отрицательным} (при $q\ge 1$). Если $\|x-y\|^2=-a$, $a>0$, то~(2) требует
\[
-(1-\varepsilon)a \le \|f(x)-f(y)\|^2 \le -(1+\varepsilon)a,
\]
то есть $-(1-\varepsilon)a \le -(1+\varepsilon)a$, что равносильно $\varepsilon\le-\varepsilon$, т.е.\ $\varepsilon\le 0$. Противоречие с $\varepsilon>0$.

\smallskip
\textbf{Следствие:} как записано, для любого источника с $q\ge1$ (где заведомо есть пары с $\|x-y\|^2<0$) \emph{не существует ни одного} $f$, удовлетворяющего~(2). Теорема~2 становится вакуумно истинной по неверной причине, а не потому, что доказано отсутствие понижения размерности.

\smallskip
\textbf{Как чинить:} использовать относительную ошибку по модулю
\[
\bigl|\,\|f(x)-f(y)\|^2-\|x-y\|^2\,\bigr|\le \varepsilon\,\bigl|\|x-y\|^2\bigr|.
\]
При таком определении рассуждение проходит (см.\ ошибку~2).
\end{critbox}

\begin{critbox}
\textbf{2. Шаг «$f$ сохраняет нулевые расстояния и только их» обоснован лишь для $\|x-y\|^2>0$.}

В тексте: «Если $\|x-y\|^2\ne0$, то $\|f(x)-f(y)\|^2\ge(1-\varepsilon)\|x-y\|^2>0$». Заключение «$>0$» верно \emph{только если} $\|x-y\|^2>0$. Для $\|x-y\|^2<0$ имеем $(1-\varepsilon)\|x-y\|^2<0$, и из этой оценки \textbf{не следует}, что образ имеет ненулевой квадрат расстояния.

Именно сохранение эквивалентности $A_{ij}=0\iff\|f(x_i)-f(x_j)\|^2=0$ (перенос реализации из источника в цель) — центральный шаг переноса; он не доказан для «отрицательных» (пространственно/времени-подобных) пар.

\smallskip
\textbf{Замечание:} при исправленном определении из ошибки~1 всё чинится: если $\|x-y\|^2=-a<0$, то $\|f(x)-f(y)\|^2\in[-(1+\varepsilon)a,\,-(1-\varepsilon)a]\subset(-\infty,0)$, т.е.\ остаётся ненулевым. Но в текущем виде это \emph{пробел в доказательстве}.
\end{critbox}

\section{Существенные замечания}

\begin{modbox}
\textbf{3. Переусиление: «матрица $A$, не реализуемая ни в каком пространстве с $p'+q'<Cn$».}

Теорема~1 гарантирует, что для \emph{каждой фиксированной} сигнатуры $s=p'+q'<Cn$ существует нереализуемая матрица — вообще говоря, \emph{своя для каждой сигнатуры}. Она не даёт напрямую \emph{единую} матрицу, плохую сразу для всех сигнатур.

Для доказательства Теоремы~2 достаточно более слабого: зафиксировав конкретное $f$ (а с ним конкретные $p',q'$), выбрать $A$, плохую именно для этой цели. С такой поправкой рассуждение корректно.

Утверждение о единой матрице тоже \emph{истинно}, но требует дополнительного объединения по $\le (Cn)^2$ сигнатурам: суммарное число реализуемых в какой-либо сигнатуре $s<Cn$ матриц не превосходит $(Cn)^2\cdot\max_s(C_0n/s)^{ns}\ll 2^{\binom n2}$. Этот шаг в тексте пропущен.
\end{modbox}

\begin{modbox}
\textbf{4. Неверная ссылка и неполно сформулированная лемма о вложении графа (Предложение~1).}

Ссылка \cite{Erdosetal} (Erdős--Frankl--Rödl, «асимптотическое число графов без фиксированного подграфа») \emph{не относится} к утверждению о вложении вершин графа в $\R^{n-1}$ с фиксированным расстоянием $\sqrt2$ на рёбрах и $\ne\sqrt2$ на не-рёбрах. Это разные темы; ссылка ошибочна.

Само используемое утверждение (существование «точного» представления графа расстоянием $\sqrt2$ в $\R^{n-1}$) правдоподобно (для полного графа — правильный симплекс), но требует \emph{корректной ссылки или явного доказательства}, включая гарантию строгого неравенства $\|y_i-y_j\|^2\ne2$ на не-рёбрах.
\end{modbox}

\section{Мелкие замечания и сомнения}

\begin{minbox}
\textbf{5. Константа $C$ не зависит от $\varepsilon$.}

Итоговая оценка $p'+q'\ge C\cdot\min(p,q)$ имеет $C$, полностью определённую подсчётом из Теоремы~1 и \emph{не зависящую от $\varepsilon$}. Гипотеза $\varepsilon$-изометрии используется лишь \emph{качественно} — только чтобы $f$ сохраняла структуру «нулевых расстояний» (световой конус). Фактически доказано более сильное и иное по духу утверждение: \emph{любое} отображение, точно сохраняющее отношение нулевого псевдорасстояния, требует линейной по $\min(p,q)$ сигнатуры цели. Это скорее результат о сохранении светового конуса, чем количественный аналог JL; аналогия с JL — поверхностная.
\end{minbox}

\begin{minbox}
\textbf{6. Роль $n=\min(p,q)$ и корректность аналогии с JL.}

Параметр $n$ одновременно есть \emph{число точек} конфигурации и привязан к сигнатуре источника ($n=\min(p,q)$). Утверждение по сути таково: «для наихудшей конфигурации из $\min(p,q)$ точек нужна цель сигнатуры, линейной по числу точек». Это законно, но в классической JL размерность цели ($O(\log N/\varepsilon^2)$) не зависит от объемлющей размерности и логарифмична по числу точек $N$; здесь же нижняя оценка линейна по числу точек. Сравнение в последнем абзаце корректно по смыслу, но формулировку стоит уточнить (что именно играет роль $N$).
\end{minbox}

\begin{minbox}
\textbf{7. Условие $m\ge D$ для оценки Милнора--Тома не проверено явно.}

Оценка $\bigl(\tfrac{C_0dm}{D}\bigr)^D$ требует $m\ge D$. В применении $m=\binom n2\approx n^2/2$, $D=ns=\alpha n^2$, так что $m\ge D\iff\alpha\lesssim1/2$ — выполнено для малого $\alpha$. Однако это условие в тексте не выписано; его стоит явно проверить.
\end{minbox}

\begin{minbox}
\textbf{8. Библиография: избыточные/неточные ссылки на количественную оценку.}

Оценку числа знаковых условий даёт \cite{BasuPollackRoy} — это верно. Ссылки \cite{Milnor1964} и \cite{Thom1965} относятся к оценкам чисел Бетти, а не непосредственно к подсчёту знаковых паттернов; их роль как источника именно \emph{этой} оценки следует уточнить (обычно это теорема типа Уоррена/Милнора--Тома в изложении Бэзу--Поллака--Роя).
\end{minbox}

\begin{minbox}
\textbf{9. Предложение~1: доказана сигнатура $(n,n-1)$, а не $(n,n)$.}

Фактически построена реализация в $\R^{n,n-1}$, затем добавлено «нулевое» измерение до $(n,n)$. Это корректно (добавление координаты, тождественно равной нулю у всех точек, не меняет расстояний), но формулировка Предложения могла бы сразу давать более сильную сигнатуру $(n,n-1)$.
\end{minbox}

\begin{minbox}
\textbf{10. $f$ квантифицируется по всему $\R^{p,q}$.}

Теорема~2 говорит о глобальных $\varepsilon$-изометриях. Поскольку реализация задействует лишь конечный набор точек, утверждение верно и в более сильной форме — для отображений, определённых лишь на конечной конфигурации. Стоит либо усилить формулировку, либо явно отметить, что глобальность не используется.
\end{minbox}

\section{Вывод}

\begin{itemize}[leftmargin=1.4em]
  \item \textbf{Раздел~2 (Предложение~1):} конструкция верна; проблема — ошибочная ссылка и неполно сформулированная вспомогательная лемма (пп.~4).
  \item \textbf{Раздел~3 (Теорема~1):} \emph{корректен} по существу; желательно явно проверить $m\ge D$ и уточнить ссылки (пп.~7--8).
  \item \textbf{Раздел~4 (Теорема~2):} содержит \textbf{критические дефекты формулировки и пробел в доказательстве} (пп.~1--2), а также переусиленное утверждение (п.~3). При исправлении определения $\varepsilon$-изометрии (относительная ошибка по модулю) и добавлении объединения по сигнатурам основной результат, по-видимому, восстанавливается, но с оговоркой, что $\varepsilon$ входит лишь качественно (п.~5).
\end{itemize}

\begin{thebibliography}{9}
\bibitem{BasuPollackRoy}
S. Basu, R. Pollack, M.-F. Roy, \emph{Algorithms in Real Algebraic Geometry}, Springer, 2006.
\bibitem{Milnor1964}
J. Milnor, \emph{On the Betti numbers of real varieties}, Proc. Amer. Math. Soc. 15 (1964), 275--280.
\bibitem{Thom1965}
R. Thom, \emph{Sur l'homologie des vari\'et\'es alg\'ebriques r\'eelles}, 1965.
\bibitem{Erdosetal}
P. Erdős, P. Frankl, V. Rődl, \emph{The asymptotic number of graphs\ldots}, Graphs and Combinatorics, 1986.
\end{thebibliography}

\end{document}
